% ============================================================================
% PART 09 - BIOLOGICAL INSIGHTS
% ============================================================================

\labelsec{biological_insights}

The three SHT units are mathematical objects, but their values encode biologically meaningful information. This chapter synthesizes findings into biologically interpretable conclusions.

\section{Sequence-Geometry-Function Chain}

A fundamental question in DNA biology is: how does nucleotide sequence encode function? The SHT framework offers a partial answer via the chain:
\begin{equation}
    \text{Sequence} \xrightarrow{\text{stacking/H-bond}} \text{Geometry}[v_0, \eta_r, \Omega] \xrightarrow{\text{SHT}} \text{Function}
\end{equation}

The validation results of Chapter~\ref{sec:unit_validation} demonstrate the first arrow. The classification results of Chapter~\ref{sec:functional_classification} demonstrate the second. Here we interpret what specific sequence patterns produce.

\subsection{A-tract Rigidity ($\rightarrow$ low Vij, low Mamton)}

A-tracts (consecutive A$\cdot$T base pairs with adenines on the same strand) are stiff because:
\begin{enumerate}
    \item Strong propeller twist of A$\cdot$T pairs forms bifurcated hydrogen bonds stabilizing the straight conformation
    \item Reduced roll angle at A$\cdot$A steps (compared to T$\cdot$A steps) minimizes bending flexibility
    \item Spine of hydration in the narrow minor groove provides additional rigidity
\end{enumerate}

\textbf{SHT signature:} $v_0 = 0.080$ (35\% below mean), $\Omega = 70.5$ (49$\times$ below control group mean). Both are consistent with straight, symmetric DNA.

\subsection{Nucleosome Geometry ($\rightarrow$ high Vij, high Mamton, moderate $\eta_r$)}

DNA wrapping around the histone octamer requires extreme local curvature. The wrapping is not uniform: DNA bends more at specific positions (SHL $\pm 1.5$, $\pm 2.5$) where the minor groove faces inward \citep{Richmond2003}. The mean curvature over the full 146-bp wrap gives $\bar{v}_0 \approx 0.30$, consistent with wrapping radius $R \approx 42$\,\AA.

The one-sided histone contact also creates strong backbone asymmetry: the ``inside'' backbone is compressed, the ``outside'' stretched. This drives $\Omega$ to extreme values ($>26{,}000$). By contrast, $\eta_r$ is not dramatically altered (1AOI: $\eta_r = 0.988$, $z = +0.54$) because the overall cross-section remains nearly circular even when asymmetrically bent.

\subsection{Protein-Induced Cross-Sectional Distortion ($\rightarrow$ low $\eta_r$)}

The most striking Shambhian outlier is 1AHD ($\eta_r = 0.710$, $z = -4.18\sigma$). This structure is a cisplatin crosslink adduct, where the platinum bridges two adjacent purines and forces a $45^\circ$ kink plus groove widening. The cross-section becomes strongly elliptic --- literally flattened --- despite only moderate curvature ($v_0 = 0.088$).

More broadly, DNA-binding proteins that induce IHF-type bending ($> 120^\circ$) or groove-specific contacts drive $\eta_r$ toward lower values by compressing the helix transversely.

\section{Nucleosome Positioning Prediction}

\subsection{The Problem}

Nucleosomes are positioned at specific genomic loci to control gene access. Current positioning predictors (NuPoP \citep{Xi2010}, NuCLE) achieve 70--80\% accuracy at specific resolution but rely on sequence features (e.g., AA/TT dinucleotide periodicity).

\subsection{SHT Contribution}

SHT offers a structural-geometry lens:
\begin{itemize}
    \item Sequences with intrinsically high curvature potential (high predicted $v_0$) are preferred nucleosome sites
    \item Stiff A-tract sequences (low predicted $v_0$) are nucleosome-depleted regions
\end{itemize}

\textbf{Preliminary evidence:} Our 33-structure labeled set shows nucleosome sequences at $v_0 > 0.25$ (both 1AOI and 1F66) while A-tract sequences cluster at $v_0 < 0.10$. A deep-learning sequence encoder trained to predict $[\bar{v}_0, \bar{\eta}_r, \bar{\Omega}]$ from sequence could then be used genome-wide to predict nucleosome affinity.

\subsection{Limitations}

The 2 nucleosome structures in our set are insufficient to train a positioning model. This requires expansion to the $\sim$50+ nucleosome PDB structures currently available.

\section{DNA Damage Detection}

\textbf{Hypothesis:} DNA damage creates characteristic geometric signatures different from canonical and protein-bound DNA.

\textbf{Evidence from 1AHD} (cisplatin adduct): $\eta_r = 0.710$ --- the most extreme Shambhian value in the non-nucleosome set. This suggests cisplatin adducts are recognizable from SHT geometry alone, without sequence knowledge.

\textbf{Predicted signatures by damage type:}
\begin{itemize}
    \item \textit{Cisplatin crosslinks}: high $\eta_r$ distortion (cross-section flattening), moderate $v_0$ increase
    \item \textit{Thymine dimers}: likely increased $v_0$ (bend at lesion site, $\sim$30$^\circ$--40$^\circ$) and moderate $\Omega$ increase
    \item \textit{Oxidative lesions (8-oxoG)}: expected small $\eta_r$ change; $v_0$ moderately affected depending on flanking sequence
\end{itemize}

These predictions are testable by running SHT on existing damage-adduct PDB structures ($> 20$ available).

\section{Chromatin Architecture}

The Vij-Mamton coupling (Chapter~\ref{sec:cross_unit_relationships}) has implications for higher-order chromatin folding. Linker DNA between nucleosomes experiences lower curvature but periodic bending at each nucleosome entry/exit. SHT applied along a full nucleosome array would yield a characteristic alternating $v_0$ profile: high spikes (nucleosome wrap regions) alternating with low valleys (linker DNA).

This signature could, in principle, be detected in chromatin fiber cryo-EM reconstructions by applying SHT to the DNA path extracted from density maps.

\section{Drug Design}

DNA-targeting drugs (actinomycin D, doxorubicin, minor-groove binders) induce geometric changes that SHT can quantify. A virtual screening pipeline could:
\begin{enumerate}
    \item Run MD simulations of drug-DNA complexes
    \item Compute SHT at each simulation frame
    \item Score changes in $[v_0, \eta_r, \Omega]$ as a geometric binding signature
    \item Select compounds that push DNA toward a target geometric state (e.g., mimic nucleosome geometry to block transcription)
\end{enumerate}

The key advantage of SHT over existing scoring functions is its direct physical interpretation: the geometric units have units and meaning, not just empirical scores.

